\section{Conclusions}

This thesis connects real wildfire detections to communication-constrained task dissemination and follow-up observation in distributed satellite systems. It contributes a traceable task pipeline, a deterministic 891-case Walker design, a temporal capacity-aware routing model, a 539-scenario replay contract, resumable Parquet output, validation before inference, and a transparent performance--resource decision method.

The earlier nominal campaign supplies a validated two-region baseline of 891 California and 891 India cases. All 1,782 cases passed completeness/configuration checks, and 17,820 independent metric comparisons passed at tolerance $10^{-6}$. It shows that regional workload and geometry change both mean outcomes and architecture rankings.

The later Fresh Full891 V2 California archive is the principal scenario evidence. It contains all 891 identifiers, 890 complete cases, 480,110 unique scenario rows after removing 4,269 resume duplicates, and 19,179 passed architecture checks. One case is incomplete and excluded from complete comparisons. These provenance details are part of the result, not administrative footnotes.

The central findings are:

\begin{enumerate}
\item Under the bounded nominal policy, mean California observation fulfilment is 61.59 percent, but strict useful-recipient completion is only 0.63 percent. Removing hop/copy limits raises these values to 77.62 and 74.69 percent at the expense of a more than tenfold increase in transferred data.
\item Central-node fraction has a non-monotonic nominal effect. Observation rises strongly between CN0 and CN10, peaks in the intermediate range, and falls at CN100; useful coverage peaks earlier. Maximum centralization is not an appropriate nominal design rule.
\item Under unlimited useful routing, performance improves and then saturates while traffic continues to grow. From CN20 to CN100, mean strict completion gains about 14.6 points while transferred data increases by roughly 211 percent.
\item Policy and topology matter. Peer and hybrid variants outperform central-only unlimited routing but require substantially more traffic; they are identical under the present recorded implementation and must not be generalized to all hybrid algorithms.
\item Larger task messages degrade observation, dissemination, and coverage. A candidate selected at 1$\times$ message size is not guaranteed to retain its position at 4$\times$.
\item Ground-station outage is the dominant tested failure, exceeding the mean losses from random satellite failure, random or targeted central-node failure, and capacity degradation. The result motivates ground-segment diversity.
\item Exact Pareto and weight-sensitivity analysis identifies \texttt{T060\_P10\_H0800\_I0986\_CN040\_F1} as the strongest balanced California compromise: 96.875 percent observation, 100 percent strict useful completion and coverage, P100 $T_{100}=66.10$ minutes, and 888.5 Mbit. It wins for performance weight 0.31--0.78 and 41.79 percent of 10,000 randomized weight draws.
\end{enumerate}

The recommended case is conditional, not universal. It assumes the California task set, fixed \SI{700}{\kilo\metre} observation radius, Munich ground station, active UHF link budget, useful-recipient semantics, unlimited diagnostic scenario for trade-space ranking, and the declared non-monetary resource weights.

\section{Answers to the Research Questions}

RQ1 is answered by the auditable conversion from Sentinel-3 fire products to multi-task, deadline-aware simulation input and the causal release--transfer--reception--observation workflow. RQ2 is answered for California priority heterogeneity and by the earlier two-region baseline, but the Fresh V2 density comparison awaits India. RQ3 is answered by the bounded non-monotonic and unlimited saturating CN-fraction responses; separate Fresh V2 first-reception latency remains unavailable. RQ4 is answered through conditional Pareto candidates and failure replay rather than a single global optimum. RQ5 is answered for task size, failure mode, and objective weights; observation-radius stability remains pending.

\section{Scientific Contribution}

The larger result is a method for deciding how much coordination capability a time-critical Earth-observation network should contain. It keeps physical opportunity, communication policy, mission fulfilment, resource burden, and robustness separate until their trade-offs are visible. The method avoids three common shortcuts: treating nominal failure labels as robustness, hiding censoring in latency means, and converting an arbitrary weighted sum into an ``optimal'' architecture without sensitivity.

Vincenzo Messina provided the communication-model and TAT-C basis. The author implemented the wildfire-data processing, integration, central-node placement, scenario execution, validation, post-processing, figures, statistical analysis, and report.

\section{Required Work before Final Submission}

The report is a thesis-quality working draft, not the final administrative submission. The remaining evidence and document actions are:

\begin{enumerate}
\item finish and validate Fresh V2 India with the same schema, completeness, duplicate, metric, and claim gates;
\item perform a small controlled observation-radius sweep on representative/Pareto cases, or document the supervisor's decision to retain it as future work;
\item extract separate first-reception latency from task-level evidence if that phrase must remain in the final Fresh V2 RQ3 answer;
\item confirm examiner, official submission date, declaration pages, acknowledgements, and faculty formatting; and
\item archive the simulation code and large raw outputs separately with immutable checksums and a credential-free release record.
\end{enumerate}

India results should be integrated as a paired Fresh V2 analysis, not appended as an isolated table. The same architecture IDs, per-task normalization, scenario-family coverage, and objective definitions should be used. If a definition changes, it should create a new evidence layer rather than overwrite California values.

\section{Outlook}

Future work can replace the central-node label with a heterogeneous subsystem model containing radio mass, power, storage, processing, reliability, and lifecycle cost. Multiple ground stations, dedicated relay shells, queues, acknowledgements, interference, simultaneous-link constraints, and fragmentable messages would make the communication model more operational. Observation scheduling could add sensor field of view, cloud/smoke probability, lighting, slew, energy, memory, and image downlink.

For wildfire applications, larger event samples and independent outcome labels could separate regional geometry from density, priority, and fire intensity. Stakeholder elicitation could replace declared engineering weights with a formally agreed value model. The present framework is structured so that these additions change one explicit layer while preserving the traceability of tasks, scenarios, and claims.
